\section{Mercy toward the neighbour and confidence before God}
The two interpretations answer different questions about the interval
before fulfilment. Chrysostom and Augustine ask what an attempted removal
of the wicked might do to the innocent and to those who can repent.
Their answer constrains the servants' remedy and demands care for the
neighbour. Jerome asks why God urges prayer for a promised gift;
Augustine asks how a sinner can hope to be heard. Their answers concern
the recipient's dependence on a mercy that has already begun to act.

The first interpretation gives special weight to fallible judgement.
An accurate account of present wrongdoing cannot guarantee knowledge of
a person's future. Nor can zeal guarantee that a remedy will spare the
good. The second gives special weight to the saving object of prayer.
The assured petitioner seeks pardon, conversion and completed salvation,
while allowing God to know the fitting gift and time. These differences
matter when a person must decide how to correct another or how to continue
praying after a desired temporal good has not arrived.

Their common ground is visible in Colossians. Those who have received
forgiveness must forgive. Those who need instruction must let Christ's
word dwell abundantly. The Secret asks mercy for the offerers themselves,
and the Postcommunion seeks the effect of a gift they have received.
The Gospel's wheat is gathered into the owner's barn; the people who
pray these prayers seek to reach that end through the Lord's mercy.
